\documentclass[11pt]{article}
\usepackage[margin=1in]{geometry}
\usepackage{graphicx}
\usepackage{booktabs}
\usepackage{amsmath}
\usepackage{hyperref}
\usepackage{natbib}

\title{Representation Instability in Induced Ontologies}

\author{Physis Research Program}
\date{\today}

\begin{document}
\maketitle

\begin{abstract}
Cross-embedder ontology agreement on induced clusterings is near chance
(ARI $\approx 0.10$). We trace the full arc from ``geometry can certify
ontology membership'' to ``geometry is unreliable for certification but
useful for proposal'' through 18 experiments on a 658-entry, 57-cell
ontology. The dominant cause is representation instability: different
embedder cones point in different directions, and no rank intervention
recovers the geometry. Anisotropy is a measured covariate (mean pairwise
cosine 0.53--0.56, slope $-0.489$ per unit), not a lever. The practical
consequence is a division of labour: geometry proposes, humans dispose.
\end{abstract}

\section{Introduction}

Ontology induction --- discovering taxonomic structure from data --- is a
founding promise of representation learning. If embeddings capture
semantic similarity, then clustering embeddings should recover ontological
structure. We test this on a 658-entry, 57-cell industrial ontology
(5 domains $\times$ 14 modes), using three vector sets over the same
entries: BGE-base with name+hints text, BGE-base with name-only text, and
Nomic-embed.

The result is a settled negative: cross-embedder adjusted Rand index (ARI)
is 0.108--0.178, near chance. The negative is now well-understood: the
disagreement has a measurable covariate (anisotropy) and a known mechanism
(cone orientation differs by embedder). But the negative does not imply
uselessness: geometry remains a powerful proposal mechanism.

\section{Experiments}

\subsection{Cross-embedder ARI}

Cross-embedder clustering ARI was measured with spherical k-means at
$k = 57$ (the number of true cells), with ARI computed against the true
cell assignments:

\begin{table}[h]
\centering
\begin{tabular}{lccc}
\toprule
\textbf{Vector set} & \textbf{Anisotropy} & \textbf{Fit ARI} & \textbf{Cross-embedder ARI} \\
\midrule
BGE name+hints & 0.558 & 0.062 & 0.108 \\
BGE name-only & 0.562 & 0.066 & 0.165 \\
Nomic-embed & 0.527 & 0.080 & 0.178 \\
\bottomrule
\end{tabular}
\caption{Anisotropy (mean pairwise cosine), fit ARI against true cells,
and cross-embedder clustering ARI. All near chance (ARI = 0).}
\label{tab:ari}
\end{table}

The covariate's predicted direction holds: the least anisotropic set
(Nomic, 0.527) has the best fit to true cells (0.080), and the slope of
fit ARI against anisotropy is $-0.489$ per unit. Three points cannot
separate a slope from noise; what is reproducible is the ARI itself.

\subsection{Injection realism}

The central confound was identified in E14: the injection policy
determines the verdict.

\begin{table}[h]
\centering
\begin{tabular}{lc}
\toprule
\textbf{Injection policy} & \textbf{AUC} \\
\midrule
Random cell (original) & 0.880 \\
Same domain, different mode & 0.827 \\
Same mode, different domain & 0.771 \\
Nearest alternative cell & 0.506 \\
Real misfilings (E16) & 0.519 \\
\bottomrule
\end{tabular}
\caption{Injection realism. Making the injected error plausible drops the
scorer to chance, and closes the 0.410/0.866 gap this track carried
unexplained throughout.}
\label{tab:injection}
\end{table}

The scorer detects \emph{distance moved}, not wrongness. A random injection
displaces an entry 0.125 below its home similarity; a plausible error,
0.029.

\subsection{Structural scorers}

Five scorers were compared under four injection policies (E15):

\begin{table}[h]
\centering
\begin{tabular}{lcc}
\toprule
\textbf{Scorer} & \textbf{Random inj.} & \textbf{Nearest inj.} \\
\midrule
COSINE & 0.882 & 0.506 \\
NAME-ONLY & 0.866 & 0.597 \\
CONSTRAINT & 0.675 & 0.612 \\
ANCHORVOC & 0.654 & 0.584 \\
ANCESTRY & 0.627 & 0.566 \\
\bottomrule
\end{tabular}
\caption{Structural scorers under two injection policies. Under random
injection, cosine leads by 0.207. Under realistic injection, cosine trails
CONSTRAINT by 0.106. Stage 12's comparison was invalid.}
\label{tab:scorers}
\end{table}

Under random injection, cosine leads. Under realistic injection, it
trails. This does not show the structural scorers work --- the best is
0.612, not a usable certifier --- but it shows the comparison that
established ``plain cosine beats every structural alternative'' was an
artifact of the injection policy.

\subsection{ANCHORVOC and the length artifact}

ANCHORVOC (lexical overlap with the cell's authored anchor hints) scored
0.664 on real misfilings (E16), beating cosine's 0.519 by +0.145. It was
the only scorer that won on real data. It was also an entry-length
artifact: entry length alone reached 0.619 of the 0.664, and no vocabulary
effect survived controlling for it (E18).

\subsection{Propose, do not certify}

The one positive that survives its construction-matched null is E21:
geometry as a proposer, not a judge. Top-3 containment of the rater's own
cell achieved 0.712 against a permuted-label null of 0.136, with the
learning curve still climbing at 436 decisions. The geometry did not get
better --- a number filed as a failure (0.447 top-1) is the same number
that makes a usable proposer once the output is three candidates instead of
one.

\section{The arc}

The research lineage is:

\begin{enumerate}
\item \textbf{Claim:} geometry can certify ontology membership.
\item \textbf{E1:} cross-embedder divergence is at chance.
\item \textbf{E2--E5:} the taxonomy is not the constraint; the grid is not
the constraint.
\item \textbf{E14:} the injection policy determines the verdict.
\item \textbf{E15:} the structural scorers do not survive the correction.
\item \textbf{E16:} all geometric scorers are at chance on real misfilings.
\item \textbf{E17--E18:} ANCHORVOC is an entry-length artifact.
\item \textbf{Verdict:} no tested mechanism detects a real misfiling. The
central negative stands, now for the right reason.
\item \textbf{E21:} geometry can propose. The task formulation changes,
not the representation.
\item \textbf{E88--E90:} anisotropy is real but not a lever; rank
intervention fails.
\end{enumerate}

Four experiments each correcting the last, ending at the position the track
held before E14 --- for entirely different and much better reasons.

\section{Discussion}

\textbf{The instrument, not the mechanism.} The cross-embedder refutation
is reproducible (a new clustering implementation lands in the same band).
But the instrument --- the corpus, the embedding model, the clustering
algorithm --- dominates the mechanism. Geometry is unreliable for
certification because the representation is unstable, not because the
ontology is absent.

\textbf{The division of labour.} The propose/certify distinction is the
finding. Geometry converts human decisions into leverage over the remaining
ones (top-3 0.712 vs.\ null 0.136). It does not replace the human. The
practical rule: sell the filter, not the judgement.

\textbf{Limits.} One corpus (658 entries, 57 cells), one embedder pair
(BGE-base vs.\ Nomic-embed), one clustering algorithm (spherical k-means).
The externalization to WordNet/SUMO/DOLCE with 3+ embedders is the
highest-priority replication. The anisotropy slope rests on three points
and we treat it as exploratory, not as an established relationship.

\section*{Reproducibility and provenance}

Every numerical claim in this paper now resolves to a stored result file.

The headline results --- the anisotropy measurements, the fit ARIs, and the
cross-embedder ARIs --- come from
\texttt{research/mapper/e88\_results.json} and
\texttt{research/mapper/e90\_results.json}, produced by
\texttt{e88\_anisotropy\_covariate.py} and \texttt{e90\_rank\_intervention.py}.

The comparison numbers from the E14--E18 chain were previously
unreproducible: those four experiments shipped only as printing scripts and
their output was never saved. They have since been re-run from the surviving
corpus and ground truth, and the recovered values are stored in
\texttt{research/e15\_e18\_rerun.json} (raw output in
\texttt{research/e15\_e18\_rerun\_raw.txt}, parsed by
\texttt{research/papers/parse\_e1518.py}). The re-run reproduced the
originally reported figures:

\begin{itemize}
\item the random-injection AUCs (COSINE 0.882, NAME-ONLY 0.866,
  CONSTRAINT 0.675, ANCHORVOC 0.654, ANCESTRY 0.627),
\item the nearest-cell injection fall (COSINE 0.506, CONSTRAINT 0.612),
\item the real-misfiling battery (COSINE 0.519, NAME-ONLY 0.508,
  CONSTRAINT 0.536, ANCHORVOC 0.664, ANCESTRY 0.534, base rate 80.1\%),
\item the length control that closes ANCHORVOC (entry length alone 0.619;
  normalised ANCHORVOC vs.\ length-only $+0.048$ [$-0.019$, $+0.116$],
  not significant).
\end{itemize}

One number deserves emphasis because the re-run changed its status from
assertion to measurement: E18's port check shows the Python scorer
reproduces the original Rust values for CONSTRAINT ($r = 0.943$) and
ANCESTRY ($r = 0.976$) but \emph{not} for COSINE ($r = 0.197$). The
divergence is expected --- the original used MiniLM, this corpus uses
BGE --- and it means the structural scorers are embedder-independent while
geometry is not. That is a stronger version of this paper's thesis and it
is now measured rather than inferred.

\end{document}
